La Marta compra 100 accions de l'empresa A, 50 de la B i 200 de la C, i paga 1.500\,€. En Pau compra
200 accions de l'empresa A, 150 de la B i 100 de la C, i paga 2.200\,€.

\begin{apartats}

\apartat[1,25]{0,75}
Planteja un sistema d'equacions i estudia amb el teorema de Rouché-Fröbenius si es pot saber el preu de
cada acció.

\begin{solucio}
Siguin $x$, $y$ i $z$ els preus, en euros, d'una acció de A, de B i de C. Dividint per 50:
\[
\left.\begin{aligned}100x+50y+200z&=1500\\200x+150y+100z&=2200\end{aligned}\right\}
\ \Longleftrightarrow\
\left.\begin{aligned}2x+y+4z&=30\\4x+3y+2z&=44\end{aligned}\right\}
\]
$\begin{vmatrix}2&1\\4&3\end{vmatrix}=2\ne0$: $\operatorname{rang}A=\operatorname{rang}A^*=2<3$. És
compatible indeterminat: hi ha infinites solucions, i amb aquestes dades no es pot saber el preu de
cada acció.
\end{solucio}

\begin{tria}{mes-informacio}
\itemtria{preus-enters}{1}{1,25}
Resol el sistema. Si se sap que els tres preus són nombres enters entre 1\,€ i 12\,€, quines
possibilitats hi ha?

\begin{solucio}
$\left(\begin{array}{ccc|c}2&1&4&30\\4&3&2&44\end{array}\right)\xrightarrow{F_2-2F_1}
\left(\begin{array}{ccc|c}2&1&4&30\\0&1&-6&-16\end{array}\right)$. Amb $z=\lambda$: $y=6\lambda-16$ i
$2x=30-y-4\lambda=46-10\lambda$, $x=23-5\lambda$.\\
Si $z$ és enter, $x$ i $y$ també. Cal $1\le6\lambda-16\le12$, és a dir, $\lambda\in\{3,4\}$, i
$1\le23-5\lambda\le12$, que tots dos compleixen. Hi ha dues possibilitats: $(x,y,z)=(8,2,3)$ i
$(x,y,z)=(3,8,4)$.
\end{solucio}

\itemtria{tercera-compra}{1}{1,25}
La Laia compra 50 accions de cada empresa i paga 750\,€. Amb aquesta dada, es pot saber el preu de cada
acció? Troba'l.

\begin{solucio}
La tercera equació és $50x+50y+50z=750$, és a dir, $x+y+z=15$. Ara
$|A|=\begin{vmatrix}2&1&4\\4&3&2\\1&1&1\end{vmatrix}=6+2+16-12-4-4=4\ne0$: el sistema és compatible
determinat, i el preu de cada acció queda determinat.\\
$\left(\begin{array}{ccc|c}1&1&1&15\\2&1&4&30\\4&3&2&44\end{array}\right)\xrightarrow[F_3-4F_1]{F_2-2F_1}
\left(\begin{array}{ccc|c}1&1&1&15\\0&-1&2&0\\0&-1&-2&-16\end{array}\right)\xrightarrow{F_3-F_2}
\left(\begin{array}{ccc|c}1&1&1&15\\0&-1&2&0\\0&0&-4&-16\end{array}\right)$.\\
$z=4$, $y=2z=8$, $x=15-8-4=3$: una acció de A val 3\,€; una de B, 8\,€, i una de C, 4\,€.
\end{solucio}
\end{tria}

\begin{nomesllarg}
\apartat{0,75}
Explica per què, amb les dades de només dues compres, no es pot saber mai el preu de les accions de tres
empreses, siguin quines siguin les quantitats i els imports.

\begin{solucio}
Dues compres donen un sistema de dues equacions amb tres incògnites. La matriu de coeficients té dues
files, i el seu rang és com a molt 2, menor que el nombre d'incògnites. Pel teorema de
Rouché-Fröbenius, el sistema no pot ser compatible determinat: o bé és incompatible (les dades serien
contradictòries), o bé és compatible indeterminat, amb infinites solucions.
\end{solucio}
\end{nomesllarg}

\end{apartats}
